% test-compat.tex --- backwards compatibility guard.
%
% This file is deliberately written in the PRE-\newagent style: it loads
% claudenotes.sty directly and uses the owner-tagged \claude family exactly as
% papers set up before agentnotes.sty existed do.  It must keep building
% unchanged, so do not modernize it --- that is the whole point of the file.
%
%   make            (or: latexmk -pdf test-compat.tex)
%
% `make compat-check` proves the guarantee rather than asserting it: it builds
% this file against the packages as they were before agentnotes.sty existed and
% compares the two renders pixel for pixel.
\documentclass{article}
\usepackage[left=1in,right=2.9in,top=1in,bottom=1in]{geometry}
\usepackage{mytodonotes}
\usepackage{claudenotes}
\setlength{\marginparwidth}{5.4cm}

\colorlet{notecolor-sam}{red!40}
\newcommand{\sam}[2][]{\note[#1]{sam}{notecolor-sam}{#2}}
\newcommand{\Sam}[2][]{\sam[inline,#1]{#2}}

\begin{document}
\section{Old-style Claude macros}

Margin note.\claude[sam]{A margin note.}  Inline note.
\Claude[sam]{An inline note.}  An author note with a Claude reply
inside it.\sam{Please check this.
\claudeResponse[sam] Checked.}

An inline suggestion: \claudeSuggest[sam]{suggested wording}, and replacement
text: \claudechange[sam]{replacement wording}.

An untagged note.\claude{No owner given.}

\begin{ClaudeSuggest}[sam]
A block suggestion, owner tagged.
\end{ClaudeSuggest}

\begin{ClaudeSuggest}
A block suggestion, untagged.
\end{ClaudeSuggest}

The three exported colors still resolve:
\textcolor{ClaudeColor}{ClaudeColor},
\textcolor{ClaudeTextColor}{ClaudeTextColor},
\textcolor{ClaudeNoteBG}{ClaudeNoteBG}.
\end{document}
